\section{All normalized scales and all centers}
\label{sec:scales}

The preceding estimate concerns a fixed scale.  To control every
$t\in[1,2]$, we track the local predicates $Q_i$ on the grids of their
own child subtrees.  The padding of each window bounds the number of
changes in these predicates as the scale varies.
The use of deterministic grid crossings to discretize scales at a fixed
center follows the same principle as
\cite[Section~3.2]{Kolountzakis1997},
\cite[Section~5]{Chlebik2015}, and
\cite[Section~3.1]{KolountzakisPapageorgiou2025}.
Here the crossings must be counted within each child subtree, so that
the bound remains compatible with the conditional trial estimate.

\begin{lemma}\label{lem:representatives}
Fix a stable center $x$ and an atom $\xi$ of the full center exposure
$\mathcal F_x$ in \eqref{eq:center-exposure} on which the first-default node $U$ exists.  Let $h$ be
its height, put $r=r_h$, and write
\[
 e_i=(U,i),\qquad b_i^*=b_{e_i}^*,\qquad 1\le i<M.
\]
There is a finite set $\mathcal R(x,U)\subset[1,2]$, depending only on
$x,U$ and the deterministic windows, with
\begin{equation}
 |\mathcal R(x,U)|\le20Mr\bigl(1+2^{2r+2}\bigr),
 \label{eq:representative-count}
\end{equation}
such that the following holds for the local predicates $Q_i$ in
\eqref{eq:local-predicate}, for every realization of the tables
consistent with $\xi$.  For each $t\in[1,2]$, some
$t'\in\mathcal R(x,U)$ satisfies
\begin{equation}
 Q_i(x+t2^{-n})=Q_i(x+t'2^{-n})
 \quad\text{for all }1\le i<M\text{ and }n\in W_{e_i}.
 \label{eq:representative-agreement}
\end{equation}
\end{lemma}

\begin{proof}
For fixed tables, $Q_i(z)$ is determined by $J_{b_i^*}(z)$.
Indeed, its initial selector has resolution $b_{e_i}\le b_i^*$,
and all selectors used to route from $U_i$ lie in that child subtree.
The incoming edge of every possible terminal leaf also has endpoint
at most $b_i^*$.  This includes the case that $U_i$ itself is a leaf,
whose terminal resolution is $b_{e_i}$, as well as leaves reached
through default children.  Since the grids are nested, the key
$J_{b_i^*}(z)$ determines every address needed for this local predicate.

For each tested pair $(i,n)$, let
\[
 \mathcal D_{i,n}(x)=
 \left\{2^n\bigl(j2^{-(b_i^*+2)}-x\bigr):j\in\mathbb Z\right\}
 \cap[1,2].
\]
These are precisely the scales at which $x+t2^{-n}$ meets a boundary
of the periodic grid $J_{b_i^*}$.  The full real-line lattice in this
definition also includes every passage through an integer, where the
periodic key returns to its first cell.

As $t$ ranges over $[1,2]$, the point traverses a closed interval of
length $2^{-n}$.  Counting lattice points in that interval gives
\[
 |\mathcal D_{i,n}(x)|\le1+2^{b_i^*+2-n}.
\]
The extra one allows both endpoints to be lattice points.  The subtree
span bound \eqref{eq:padding} gives $b_i^*-a_{e_i}+1\le2r$; since
$n\ge a_{e_i}$,
\[
 |\mathcal D_{i,n}(x)|\le1+2^{2r+2}.
\]

Take the union of these sets and the two endpoints:
\[
 \mathcal D(x,U)=\{1,2\}\cup
 \bigcup_{1\le i<M}\ \bigcup_{n\in W_{e_i}}\mathcal D_{i,n}(x).
\]
List its distinct elements as $d_1=1<\cdots<d_L=2$, and set
\[
 \mathcal R(x,U)=\mathcal D(x,U)
 \cup\left\{\frac{d_j+d_{j+1}}2:1\le j<L\right\}.
\]
There are $(M-1)r$ tested pairs, so
\begin{align*}
 L&\le2+(M-1)r\bigl(1+2^{2r+2}\bigr),\\
 |\mathcal R(x,U)|=2L-1
 &\le2(M-1)r\bigl(1+2^{2r+2}\bigr)+3\\
 &\le20Mr\bigl(1+2^{2r+2}\bigr).
\end{align*}
On each open interval $(d_j,d_{j+1})$, every relevant grid key is
constant, hence so is every tested local predicate.  Its midpoint
therefore represents the entire interval.  Each breakpoint and each
endpoint is included separately and represents its own value.  This
proves \eqref{eq:representative-agreement}, including the values at
the boundaries of the half-open cells.
\end{proof}

The representatives use only the local predicates.  A point for which
$Q_i$ vanishes may take its actual route through another subtree with
a finer grid.  That route does not enter the count: we need only the
implication from a successful local predicate to a hit of $B$.

\begin{proposition}\label{prop:raw-failure}
For every fixed stable center $x$,
\begin{equation}
 \Pbb\bigl(\exists t\in[1,2]\ \forall n\in\mathcal N:
 x+t2^{-n}\notin B\bigr)\le2p.
 \label{eq:raw-failure}
\end{equation}
\end{proposition}

\begin{proof}
Condition first on an atom $\xi$ of $\mathcal F_x$ in \eqref{eq:center-exposure} on which the
first-default node $U$ exists, and retain the notation of
Lemma~\ref{lem:representatives}.  The node $U$ and the entire set
$\mathcal R(x,U)$ are fixed on this atom, before any remaining
selector or terminal label is revealed.  Thus Lemma~\ref{lem:fixed-scale}
applies separately to every $t'\in\mathcal R(x,U)$ and gives
\[
 \Pbb\bigl(Q_i(x+t'2^{-n})=0\text{ for all tested }(i,n)
       \mid\xi\bigr)
 =(1-p/2)^{(M-1)r}
 \le\exp\!\left(-\frac{p(M-1)}2r\right).
\]

If some scale $t$ misses $B$ at every index in $\mathcal N$, then
Lemma~\ref{lem:routing} implies that all its tested local predicates
vanish.  By Lemma~\ref{lem:representatives}, they all vanish at one
of the representative scales as well.  The conditional union bound
and \eqref{eq:entropy-budget} therefore give
\begin{align*}
 &\Pbb\bigl(\exists t\in[1,2]\ \forall n\in\mathcal N:
 x+t2^{-n}\notin B\mid\xi\bigr)\\
 &\qquad\le
 20Mr\bigl(1+2^{2r+2}\bigr)
 \exp\!\left(-\frac{p(M-1)}2r\right)<p.
\end{align*}
No independence between different representative scales is required.

This bound holds on every center-exposure atom that has a first
default, since every corresponding window length is at least $r_*$.
Averaging over those atoms introduces no factor for the number of
possible nodes.  By Lemma~\ref{lem:first-default}, the remaining
event, on which the center's path never chooses the default child,
has probability less than $p$.  Bounding failure by one on that
event yields \eqref{eq:raw-failure}.
\end{proof}
